% !TEX root = ../bb_AiRa_Kappaleet.tex

%% HARRIS: Chapter 10

% ö = \%c3\%b6
% ä = \%C3\%A4

\xdef\sectiontitle{\IIsectiontitle} %aiheuttama
\TOCrefs

%%%%%%%%%%%%%%%
\begin{frame}[label=2_2_a]%[label=2_0_TOC1]
\frametitle{\texorpdfstring{\culine[\sectiontitleulinecolor]{\sectiontitle}}{\sectiontitle} }
\label{\TOCname}

\vspace{\chapterTOCvksip}
\begin{itemize}[leftmargin=0.5cm,itemsep=3mm]
    \item[2.1]<1-> \hyperlink{2_1_a}{\IIaaSubsectiontitle}
    \item[2.2]<1-> \hyperlink{2_2_a}{\CDAlert[red]{\IIaSubsectiontitle}}
    \item[2.3]<1-> \hyperlink{2_3_a}{\IIbSubsectiontitle}	
    \item[2.4]<1-> \hyperlink{2_4_a}{\IIcSubsectiontitle}	
    \item[2.5]<1-> \hyperlink{2_5_a}{\IIdSubsectiontitle}
    \item[2.6]<1-> \hyperlink{2_6_a}{\IIeSubsectiontitle}
    \item[2.7]<1-> \hyperlink{2_7_a}{\IIfSubsectiontitle}
    \item[2.8]<1-> \hyperlink{2_8_a}{\IIgSubsectiontitle}
    \item[2.9]<1-> \hyperlink{2_9_a}{\IIhSubsectiontitle}
    \item[2.10]<1-> \hyperlink{2_10_a}{\IIiSubsectiontitle}
\end{itemize}

%\endofslide 
\end{frame}
%%%%%%%%%%%%%%%

\xdef\sectiontitle{\IIaSubsectiontitle} 

%%%%%%%%%%%%%%%
\begin{frame}
\frametitle{\texorpdfstring{\culine[\sectiontitleulinecolor]{\sectiontitle}}{\sectiontitle} }
\framesubtitle{Atomien välisten sidosten muodostuminen}
\appear{}

	% 6.5 cm = midway, 10 cm = 3/4 
\vspace{1mm}
\begin{minipage}{8cm}
\begin{itemize}
	\item Kun on \Alert[fuchsia]{energeettisesti edullista}\appear{, atomit muodostavat molekyylejä}\appear{, eivätkä pysy atomaarisena kaasuna.} \appear{Molekyylin muodostuessa} \appear{sen \CDAlert[red]{valenssielektronit} järjestäytyvät uudelleen} \appear{minimoidakseen systeemin kokonaisenergian}
	\end{itemize}
\end{minipage}

\vspace{3mm}

\begin{itemize}
\item \CDAlert[red]{$\mathrm{H}_{2}{ }^{+}$-molekyyli}:
\item[] \begin{itemize}
	\item Yksinkertainen esimerkki on $\mathrm{H}_{2}{ }^{+}$-molekyyli\appear{, joka on kahdesta protonista} \appear{ja yhdestä elektronista koostuva ioni}
	\appear{\xdef\loc{north east}%
\begin{tikzpicture}[remember picture,overlay] % anchor=south
    \node (A) [yshift=\figYshift,xshift=\figXshift, anchor=\loc] at (current page.\loc) {\includegraphics[page=1,width=5.5cm]{Figures_booklet/SOM_10_Bonding_figures/SOM_Bonding_Figure_3.png}}; %scale=\figscale. % 8cm max height
\end{tikzpicture}\footnoteleft{\HarrisCR}}%
	\item Kokonaisenergia muodostuu:
\begin{itemize}
	\item Elektronin energia \appear{(kineettinen} \appear{+ negatiivinen sähköstaattinen (attraktiivinen) potentiaalienergia)}
\item \CDAlert[blue]{Sähköinen} hylkivä voima kahden protonin välillä
\end{itemize}

\end{itemize}

\end{itemize}

\endofslide 
%\endofsection
\end{frame}
%%%%%%%%%%%%%%%

%%%%%%%%%%%%%%%
\begin{frame}
\frametitle{\texorpdfstring{\culine[\sectiontitleulinecolor]{\sectiontitle}}{\sectiontitle} }
\framesubtitle{Atomien välisten sidosten muodostuminen}
\appear{}

\begin{itemize}
	\item \CDAlert[red]{Elektroni} $\mathrm{H}_{2}{ }^{+}$-molekyylissä:
	\item[] \begin{itemize} \small
	\item Jos protonien välinen etäisyys olisi suuri\appear{, niin elektronit olisivat sidoksissa jompaankumpaan protoniin}\appear{, jolloin niiden \Alert[red]{sähköinen potentiaalienergia} olisi $-13,6~\mathrm{eV}$} 
\item Toisessa ääripäässä\appear{, yhteen liittyneet protonit}\appear{, näyttäisivät He atomilta}\appear{, jolloin sähköinen potentiaalienergia olisi  $-54,4~\mathrm{eV}$}
\implies{Oletamme elektronin olevan alimmalla energiatilalla\appear{, mikä riippuu protonien välisestä etäisyydestä} \appear{(energia pienenee etäisyyden pienentyessä)} }
\end{itemize}


\item \CDAlert[blue]{Protonit} $\mathrm{H}_{2}{ }^{+}$-molekyylissä:
\item[] \begin{itemize} \small
	\item Positiivinen kontribuutio \appear{(\Alert[blue]{hylkivä})} \appear{kokonaisenergiaan kasvaa etäisyyden pienentyessä}\appear{, dramaattisesti}\appear{, kun etäisyydestä tulee todella pieni}
\item Kaukana olevilla protoneilla\appear{ vaikutus on mitätön}
\end{itemize}

\item \CDAlert[fuchsia]{Yhteensä}
\item[] \begin{itemize} \small
	\item Molekyylejä kuten $\mathrm{H}_{2}{ }^{+}$ on olemassa\appear{, eli täytyy löytyä \CDAlert[fuchsia]{potentiaalienergian minimi} jollain protonien välisellä etäisyydellä}
\end{itemize}


\end{itemize}

\endofslide 
%\endofsection
\end{frame}
%%%%%%%%%%%%%%%

%%%%%%%%%%%%%%%
\begin{frame}
\frametitle{\texorpdfstring{\culine[\sectiontitleulinecolor]{\sectiontitle}}{\sectiontitle} }
\framesubtitle{Atomien välinen potentiaalienergia}
%\vspace{-1.5\topsep}%\topsep
%\vspace{-6mm}
\appear{\xdef\loc{north east}%
\begin{tikzpicture}[remember picture,overlay]% anchor=south
    \node (A) [yshift=\figYshift,xshift=\figXshift, anchor=\loc] at (current page.\loc) {\includegraphics[page=1,width=5.5cm]{Figures_booklet/SOM_10_Bonding_figures/SOM_Bonding_Figure_3.png}};%scale=\figscale. % 8cm max height
\end{tikzpicture}%
\footnoteleft{\HarrisCR}}%
% 6.5 cm = midway, 10 cm = 3/4 
\begin{minipage}{8cm}
\begin{itemize}[itemsep=5mm]
	\item Tiedämme elektronin \CDAlert[red]{potentiaalienergian} \appear{kahden protonin sähköstaattisessa kentässä:}
\[ 
 U(\mathbf{r})  \equal \frac{1}{4 \pi \varepsilon_{0}} \appear{\LP} \frac{-e^{2}}{| \mathbf{r}-\Frac{1}{2} a \hat{\mathbf{x}} |}  \plus \frac{-e^{2}}{ | \mathbf{r}+\Frac{1}{2} a \hat{\mathbf{x}} |} \;\appear{\RP} 
 \]
\item Kun tämä potentiaalienergia sijoitetaan 
\end{itemize}
\end{minipage}
%\vspace{-2.3mm}
\begin{itemize}
\item[] \CDAlert[blue]{Schrödingerin yhtälöön}\appear{, yhtälöä ei valitettavasti voida enää ratkaista analyyttisesti}

\item \CDAlert[red]{Numeerinen ratkaisu} kertoo meille että \appear{ elektronin potentiaalienergian minimi} \appear{(protonien etäisyyden funktiona)} \appear{on \CDAlert[red]{$-16,3 \mathrm{\;eV}$}}\appear{, ja saadaan \CDAlert[red]{etäisyydellä $a_{\min }=0,11 \mathrm{~nm}$}}\appear{, mikä on vain noin \CDAlert[tunired]{kaksi kertaa Bohrin säde}}
\end{itemize}

\endofslide 
%\endofsection
\end{frame}
%%%%%%%%%%%%%%%

%%%%%%%%%%%%%%%
\begin{frame}
\frametitle{\texorpdfstring{\culine[\sectiontitleulinecolor]{\sectiontitle}}{\sectiontitle} }
\framesubtitle{Aaltofunktioiden parillisuus} 
\appear{}
\vspace{-1.5\topsep}%\topsep
\xdef\loc{south east}%
\begin{tikzpicture}[remember picture,overlay] % anchor=south
    \node (A) [yshift=-0.25*\figYshift,xshift=\figXshift, anchor=\loc] at (current page.\loc) {\includegraphics[page=1,height=6cm]{Figures_booklet/SOM_10_Bonding_figures/SOM_Bonding_Figure_4.png}}; %scale=\figscale. % 8cm max height
\end{tikzpicture}%
\footnote{\HarrisCR}%
\begin{itemize}
	\item Molekyylillä on kaksi 1s-orbitaalilla olevaa elektronia \appear{joiden aaltofunktiot ovat $\Psi_{1 s} \textrm{((}r_{1} \textrm{))}$} \appear{ja $\Psi_{1 s} \textrm{((}r_{2} \textrm{))}$}

\item Kun protonit lähestyvät toisiaan\appear{, meidän täytyy laskea aaltofunktioiden \CDAlert[fuchsia]{superpositio}:}
\end{itemize}

% 6.5 cm = midway, 10 cm = 3/4 
\begin{minipage}{6cm}
\begin{itemize}
	\item[] \appear{Schrödingerin aaltoyhtälö antaa kahdentyyppisiä ratkaisuja}\appear{, \CDAlert[red]{symmetrisiä} \CDAlert[red]{(parillisia)}:}
\[ 
 \Psi_{\text {symmetrinen}}  \equal \frac{1}{\sqrt{2}} \appear{\textrm{((} \Psi_{100} (r_{1} )}  \plus \appear{\Psi_{100} (r_{2} ) \textrm{))}} 
 \]
\item[] ja \CDAlert[blue]{antisymmetrisiä (parittomia)}:
\[ 
 \Psi_{\text {antisymmetrinen}}  \equal \frac{1}{\sqrt{2}} \appear{\textrm{((}\Psi_{100} (r_{1} )}  \minus \appear{\Psi_{100} (r_{2} ) \textrm{))}} 
 \]
\end{itemize}
\end{minipage}

\endofslide 
%\endofsection
\end{frame}
%%%%%%%%%%%%%%%

%%%%%%%%%%%%%%%
\begin{frame}
\frametitle{\texorpdfstring{\culine[\sectiontitleulinecolor]{\sectiontitle}}{\sectiontitle} }
\framesubtitle{Aaltofunktioiden parillisuus}
\appear{}
\vspace{-1.5\topsep}%\topsep
\xdef\loc{south east}%
\begin{tikzpicture}[remember picture,overlay]% anchor=south
    \node (A) [yshift=-0.25*\figYshift,xshift=\figXshift, anchor=\loc] at (current page.\loc) {\includegraphics[page=1,height=6cm]{Figures_booklet/SOM_10_Bonding_figures/SOM_Bonding_Figure_4.png}};%scale=\figscale. % 8cm max height
\end{tikzpicture}%
\footnote{\HarrisCR}%
\begin{itemize}[topsep=-1cm]
	\item Molekyylin \CDAlert[red]{perustilaa} vastaava aaltoyhtälö \CDAlert[red]{on symmetrinen}\appear{, jolloin suurin osa sen todennäköisyysjakaumasta}\appear{/varaustiheydestä} \appear{on protonien välissä olevassa tilassa}

\item \CDAlert[tunired]{Vetää molempia puoleensa}\\ 
	\CDAlert[tunired]{protonit jakavat} elektroni-\\
	pilven\appear{, mistä molekyylin \\
	 sidosenergian minimoituminen \\ johtuu}
\end{itemize}

%\xdef\loc{south east}
%\begin{tikzpicture}[remember picture,overlay] % anchor=south
%    \node (A) [yshift=-0.25*\figYshift,xshift=\figXshift, anchor=\loc] at (current page.\loc) {\includegraphics[page=1,height=6cm]{Figures_booklet/SOM_10_Bonding_figures/SOM_Bonding_Figure_4.png}}; %scale=\figscale. % 8cm max height
%\end{tikzpicture}
%\footnoteleft{\HarrisCR}

\endofslide 
%\endofsection
\end{frame}
%%%%%%%%%%%%%%%

%%%%%%%%%%%%%%%
\begin{frame}
\frametitle{\texorpdfstring{\culine[\sectiontitleulinecolor]{\sectiontitle}}{\sectiontitle} }
\framesubtitle{Kovalenttinen sidos}
\appear{}

\begin{itemize}
	\item $\mathrm{H}_{2}{ }^{+}$-ionin koostumusta \appear{voidaan muuttaa lisäämällä toinen elektroni}\appear{, jolloin \CDAlert[red]{syntyy sähköisesti neutraali $\mathrm{H}_{2}$ molekyyli} }
	
	\item Silloin olisi kaksi elektronia\appear{ perustilassaan}\appear{, molemmat samassa alimmassa tilassa} 
	\implies{ \appear{\CDAlert[red]{Tämä energian minimoituminen jakamalla valenssielektroneita} \CDAlert[red]{joilla on vastakkaiset spinit}} \appear{tunnetaan  \CDAlert[red]{kovalenttisena sidoksena}} }

\item Elektronipari on todennäköisimmin protonien välissä \appear{ja vetää molempia ioneja puoleensa} \appear{(nyt protoneita).} \appear{Molemmat atomit (protonit) $\mathrm{H}_{2}$:ssa} \appear{'näkevät' molemmat elektronit}\appear{, jolloin syntyy erittäin stabiili rakenne}

\item \appear{\CDAlert[blue]{Sama tapahtuu valenssin $-1$ alkuaineilla}}\appear{, samanlaisia elektronin jakavia sidoksia muodostuu myös $\mathrm{F}_{2}$-}\appear{, $\mathrm{Cl}_{2}$-}\appear{, $\mathrm{Br}_{2}$-} \appear{ja $\mathrm{I}_{2}$-molekyyleille}

\item $\mathrm{O}_{2}$\appear{, atomit \CDAlert[tunired]{jakavat neljä elektronia}}\appear{, $\mathrm{N}_{2}$ kuusi elektronia jaetaan}
\end{itemize}

\endofslide 
%\endofsection
\end{frame}
%%%%%%%%%%%%%%%

%%%%%%%%%%%%%%%
\begin{frame}
\frametitle{\texorpdfstring{\culine[\sectiontitleulinecolor]{\sectiontitle}}{\sectiontitle} }
\framesubtitle{Sitovat ja hajottavat sidokset}
\appear{}
%\vspace{-1\topsep}%\topsep
\vspace{-6mm}
\footnote{\HarrisCR}
\xdef\loc{south east}
\begin{tikzpicture}[remember picture,overlay] % anchor=south
    \node (A) [yshift=-0.25*\figXshift,xshift=\figXshift, anchor=\loc] at (current page.\loc) {\includegraphics[page=1,height=6.7cm]{Figures_booklet/SOM_10_Bonding_figures/SOM_Bonding_Figure_4.png}}; %scale=\figscale. % 8cm max height
\end{tikzpicture}%
\begin{itemize}[itemsep=1.5mm] \semismall
	\item Atomit näyttävät muodostavan jalokaasujen tapaisia yhdisteitä \appear{jakamalla valenssielektroneitaan.} \appear{Tämä on kuitenkin yksinkertaistettu kuvaus}\appear{, koska aaltofunktiot näyttävät erilaisilta kuin jalokaasujen tapauksessa}
	\item Myös '\CDAlert[fuchsia]{heikosti sitoutuvia} elektroneita' voidaan jakaa
	 % 6.5 cm = midway, 10 cm = 3/4 
\end{itemize}
\vspace{3mm}
\begin{minipage}{5.5cm} 
\begin{itemize}[itemsep=1.5mm] \semismall

\item Muistellaan $\mathrm{H}_{2}{ }^{+}$-ionin aaltofunktiota $\Psi$: 

\item[] \begin{itemize}[leftmargin=0.2cm,rightmargin=0.0cm,itemsep=1.8mm] \small
	\item \CDAlert[red]{Parillisen} ratkaisun energia on matalampi\appear{, sillä molemmat protonit ovat lähellä keskellä olevaa jaettua elektronipilveä.} \appear{Tämä kon\-fi\-guraatio tunnetaan nimellä \CDAlert[red]{sitovan sigmasidoksen orbitaali}}\appear{, $\sigma 1 s$.}
	 
	\item \CDAlert[blue]{Pariton} ratkaisu tunnetaan nimellä \Alert[blue]{hajottavan sigmasidoksen orbi\-taalina}\appear{, $\sigma^{*} 1 s$.} \appear{Sen energia on suurempi} \appear{koska aaltofunktion keskellä ei ole elektroneita (nollakohta).}
\end{itemize}
\end{itemize}
\end{minipage}

\endofslide 
%\endofsection
\end{frame}
%%%%%%%%%%%%%%%

%%%%%%%%%%%%%%%
\begin{frame}
\frametitle{\texorpdfstring{\culine[\sectiontitleulinecolor]{\sectiontitle}}{\sectiontitle} }
\framesubtitle{Sitovat ja hajottavat sidokset}
\appear{}
%\vspace{-1.5\topsep}%\topsep
\footnote{\HarrisCRshort}
\xdef\loc{north east}%
\begin{tikzpicture}[remember picture,overlay] % anchor=south
    \node (A) [yshift=0.5*\figYshift,xshift=\figXshift, anchor=\loc] at (current page.\loc) {\includegraphics[page=1,width=4.0cm]{Figures_booklet/SOM_10_Bonding_figures/SOM_Bonding_Figure_5.png}}; %scale=\figscale. % 8cm max height
\end{tikzpicture}%
\xdef\loc{south east}%
\begin{tikzpicture}[remember picture,overlay] % anchor=south
    \node (A) [yshift=-1.25*\figYshift,xshift=\figXshift, anchor=\loc] at (current page.\loc) {\includegraphics[page=1,width=4.0cm]{Figures_booklet/SOM_10_Bonding_figures/SOM_Bonding_Figure_6.png}}; %scale=\figscale. % 8cm max height
\end{tikzpicture}%
% 6.5 cm = midway, 10 cm = 3/4 
\begin{minipage}{9.5cm}
\begin{itemize}[itemsep=2.0mm] \seminormal
	\item Kuvassa 5 verrataan atomitilojen \CDAlert[fuchsia]{orbitaaleja}:\\ \appear{$1 s$ sitova sigmasidos} \appear{ja $1 s$ hajottava sigmasidos}

\item Kuvassa $\pm$-merkit eivät kuvasta sähkövarausta\appear{, vaan aaltofunktion vaihetta} \appear{(millä ei ole merkitystä todennäköisyyksiä tarkastellessa)}

\item \appear{\CDAlert[red]{Parillinen $\sigma 1s$} johtaa} \appear{$\Psi$:ien \Alert[red]{samanvaihei\-seen/konstruktiiviseen} yhdistymiseen}

\item \appear{\CDAlert[blue]{Parittomassa $\sigma^{*} 1s$}} \appear{$\Psi$:t ovat \Alert[blue]{vastakkaisissa vaiheissa}}

\item Tilojen energiakaavio näyttää molemmat tilanteet\appear{, kun protonit ovat kaukana toisistaan} \appear{ja kun ne ovat lähekkäin.} \appear{Käyttäytyminen on samanlaista kuin $N$:n neliökaivon tapauksessa} 
\end{itemize}
\end{minipage}

\vspace{2mm}
\begin{itemize}[itemsep=2.0mm]
\item \Alert[fuchsia]{Energiatilat loittonevat toisistaan atomien etäisyyden pienentyessä}
\end{itemize}

\endofslide 
%\endofsection
\end{frame}
%%%%%%%%%%%%%%%

%%%%%%%%%%%%%%%
\begin{frame}
\frametitle{\texorpdfstring{\culine[\sectiontitleulinecolor]{\sectiontitle}}{\sectiontitle} }
\framesubtitle{Sitovat ja hajottavat 1s-orbitaalit}
\appear{}

\begin{itemize}
	\item Käsiteltyjen molekyylien tapauksessa \appear{energiatasojen} \appear{(ja spinien)} \appear{hahmottelu havainnollistaa tilannetta: \\[2mm]}%
	\appear{\xdef\loc{south}%
\begin{tikzpicture}[remember picture,overlay] % anchor=south
    \node (A) [yshift=-0.25*\figYshift,xshift=0*\figXshift, anchor=\loc] at (current page.\loc) {\includegraphics[page=1,width=11cm]{Figures_booklet/SOM_10_Bonding_figures/SOM_Bonding_Figure_7.png}}; %scale=\figscale. % 8cm max height
\end{tikzpicture}\footnoteleft{\HarrisCR}}%
 \appear{$\mathrm{H}_{2}{ }^{+}$}\appear{, $\mathrm{H}_{2}$}\appear{, $\mathrm{He}_{2}{ }^{+}$} \appear{ja $\mathrm{He}_{2}$} \appear{(jota ei löydy luonnosta)}

\item Hajottava orbitaali on korkeammalla energialla \appear{kuin sitova $1 s$ orbitaali}

\item $\mathrm{He}_{2}{ }^{+}$-ioni voi muodostaa molekyylin\appear{, koska kaksi elektronia $\sigma 1 s$ tilalla} \appear{ja yksi $\sigma^{*} 1s$ tilassa} \appear{johtaa \Alert[fuchsia]{kokonaisenergian minimoitumiseen verrattuna kolmeen atomaariseen $1s$ tilaan (yksittäinen He ja He$^+$)}}
\end{itemize}


\endofslide 
%\endofsection
\end{frame}
%%%%%%%%%%%%%%%

%%%%%%%%%%%%%%%
\begin{frame}
\frametitle{\texorpdfstring{\culine[\sectiontitleulinecolor]{\sectiontitle}}{\sectiontitle} }
\framesubtitle{Sitovat ja hajottavat 1s-orbitaalit}
\appear{}

	% 6.5 cm = midway, 10 cm = 3/4 
\begin{minipage}{8.7cm}
\begin{itemize}[itemsep=1.6mm] \seminormal
	\item Pallomaiset $2 s$ atomiorbitaalit/-tilat käyttäytyvät kuin \appear{$1 s$ orbitaalit} 
	
	\item \CDAlert[red]{$2p$-orbitaalit} käyttäytyvät monimutkaisemmin \appear{(\CDAlert[red]{jännempiä} sidoksia)}%, as we will see next

\item Kuvassa 8 näkyy \appear{kolme $2 p$-orbitaalia}\appear{, koh\-ti\-suorassa toisiinsa nähden}
	\appear{\xdef\loc{north east}%
\begin{tikzpicture}[remember picture,overlay]% anchor=south
    \node (A) [yshift=\figYshift,xshift=\figXshift, anchor=\loc] at (current page.\loc) {\includegraphics[page=1,width=5cm]{Figures_booklet/SOM_10_Bonding_figures/SOM_Bonding_Figure_8.png}};%scale=\figscale. % 8cm max height
\end{tikzpicture}%
\footnote{\HarrisCR}%
}

\item Kulmariippuvuudet kolmelle \Alert[tunired]{samaenergiselle} vetymäiselle tilalle ovat:
\[ 
 \psi_{2,1,0} \propto \appear{\cos \theta} \qquad \appear{ja} \qquad \appear{\psi_{2,1, \pm 1}} \propto \appear{\sin \theta} \,\appear{e^{\pm i \phi}} 
 \]
 \end{itemize}
\end{minipage}
%\vspace{3mm}

\begin{itemize}[itemsep=1.5mm]
\item Jotka voidaan yhdistää \CDAlert[red]{kolmeksi ekvivalentiksi tilaksi}:
\[
\begin{aligned}
\appear{\psi_{2 p_{z}} &=\psi_{2,1,0}}  &\propto& \appear{\;\cos \theta}  \\
 \appear{\psi_{2 p_{x}}} & \equal  \appear{\psi_{2,1,+1}}  \plus \appear{\psi_{2,1,-1}}  &\propto& \;\appear{\sin \theta \cos \phi} \\ 
 \appear{\psi_{2 p_{y}}&=}\appear{\psi_{2,1,+1}}  \minus \appear{\psi_{2,1,-1}} &\propto& \;\appear{\sin \theta \sin \phi} \qquad
\end{aligned}
\]

\end{itemize}


\endofslide 
%\endofsection
\end{frame}
%%%%%%%%%%%%%%%

%%%%%%%%%%%%%%%
\begin{frame}
\frametitle{\texorpdfstring{\culine[\sectiontitleulinecolor]{\sectiontitle}}{\sectiontitle} }
\framesubtitle{2p-orbitaalit ja $\pi$-sidokset}
\appear{}

% 6.5 cm = midway, 10 cm = 3/4 
\begin{minipage}{9.5cm}
	\begin{itemize}[itemsep=3mm]
	\item Kahdella identtisellä atomilla \appear{on \Alert[fuchsia]{kuusi erilaista} $2 p$-orbitaalia}

\item Tilojen \CDAlert[blue]{$2 p_{x}$}\appear{ varaustiheys on suurin molekyylin akselilla}\appear{, ja ne tunnetaan nimellä \CDAlert[blue]{$\sigma$-sidos}}

\item On olemassa sekä sitovia $\sigma 2 p_{x}$ \appear{että että hajottavia $\sigma^{*} 2 p_{x}$ tiloja}

\item Muilla suunnilla on vastaavasti $2 p_{y}$ ja $2 p_{z}$  tilat. \appear{Näitä kutsutaan \CDAlert[red]{$\pi$-sidoksiksi}} 

\item Tilat ovat: \appear{$\pi 2 p_{y}$}\appear{, $\pi 2 p_{z}$}\appear{, $\pi^{*} 2 p_{y}$} \appear{ja $\pi^{*} 2 p_{z}$}
\end{itemize}
\end{minipage}
\vspace{3mm}

\xdef\loc{east}
\begin{tikzpicture}[remember picture,overlay] % anchor=south
    \node (A) [yshift=0*\figYshift,xshift=\figXshift, anchor=\loc] at (current page.\loc) {\includegraphics[page=1,height=8.5cm]{Figures_booklet/SOM_10_Bonding_figures/SOM_Bonding_Figure_9.png}}; %scale=\figscale. % 8cm max height
\end{tikzpicture}
\footnote{\HarrisCR[1-]}

\endofslide 
%\endofsection
\end{frame}
%%%%%%%%%%%%%%%

%%%%%%%%%%%%%%%
\begin{frame}
\frametitle{\texorpdfstring{\culine[\sectiontitleulinecolor]{\sectiontitle}}{\sectiontitle} }
\framesubtitle{2p-orbitaalit ja $\pi$-sidokset}
\appear{}

\begin{itemize}
	\item Mielenkiintoisesti kun \CDAlert[blue]{samaytimiset} \CDAlert[tuniblue]{kaksiatomiset} molekyylit muodostavat molekyylejä näiden sidoksien avulla\appear{, energiatasot $\sigma$- ja $\pi$-tiloilla eivät \CDAlert[fuchsia]{välttämättä} \Alert[fuchsia]{ole samassa järjestyksessä eri molekyyleille}}
\end{itemize}
\appear{\xdef\loc{south east}%
\begin{tikzpicture}[remember picture,overlay] % anchor=south
    \node (A) [yshift=-0.35*\figYshift,xshift=\figXshift, anchor=\loc] at (current page.\loc) {\includegraphics[page=1,width=7.5cm]{Figures_booklet/SOM_10_Bonding_figures/SOM_Bonding_Figure_10.png}};%scale=\figscale. % 8cm max height
\end{tikzpicture}\footnote{\HarrisCR}}%
% 6.5 cm = midway, 10 cm = 3/4 
\begin{minipage}{6.25cm}
\begin{itemize}
	\item Esim.~$\mathrm{N}_{2}, \mathrm{O}_{2}$ ja $\mathrm{F}_{2}$:
	
\begin{itemize}[itemsep=2.2mm]
\item $1 s$ elektronit ovat sitoutuneita vastaaviin atomeihinsa

\item  \Alert[fuchsia]{Energiatasojen järjestys vaihtelee}\appear{, eli \CDAlert[blue]{$\sigma 2 p$} on kor\-keam\-malla kuin \CDAlert[red]{$\pi 2 p$} $\mathrm{N}_{2}$:lle}\appear{, mutta alempana kuin  $\mathrm{O}_{2}$:lla ja $\mathrm{F}_{2}$:lla}

\item Kaksi $\pi^{*} 2 p$ elektronia $\mathrm{O}_{2}$:ssa ovat erilaisilla tiloilla \appear{($y$ vs. $z$)}
\end{itemize}
\end{itemize}
\end{minipage}
\vspace{3mm}



\endofslide 
%\endofsection
\end{frame}
%%%%%%%%%%%%%%%

%%%%%%%%%%%%%%%
\begin{frame}
\frametitle{\texorpdfstring{\culine[\sectiontitleulinecolor]{\sectiontitle}}{\sectiontitle} }
\framesubtitle{Polaarinen kovalenttinen sidos}
\appear{}

\begin{itemize}
	\item Jos alkuaineet \CDAlert[red]{kaksiatomisessa molekyylissä} ovat erilaiset\appear{, niin elekronit eivät jakaannu tasaisesti.} \appear{Tämä johtaa \CDAlert[red]{polaariseen kovalenttiseen sidokseen}}
\end{itemize}
\vspace{3mm}

% 6.5 cm = midway, 10 cm = 3/4 
\begin{minipage}{8.5cm}
\begin{itemize}
	\item Esim.~$\mathrm{HF}$:ssa\appear{, vedyn $\mathrm{n}=1$} \appear{ja fluorin $\mathrm{n}=2$ elektronit} \appear{jakaantuvat epäsymmetrisesti}

\item Epäsymmetria on erityisen suurta kun toisen atomin  \appear{\CDAlert[red]{valenssi on $+1$}} \appear{(haluaa yhden elektronin)} \appear{ja toisen} \appear{\CDAlert[blue]{valenssi on $-1$}} \appear{(haluaa luovuttaa yhden elektronin)}

\implies{ Atomit  \CDAlert[fuchsia]{eivät jaa elektroneitaan}\appear{, vaan toinen \\
		atomi luovuttaa 'ylimääräisen' elektronin  \\ toiselle} } 
\implies{ Tämä johtaa \CDAlert[fuchsia]{ionisidoksen} muodostumiseen }
%In practice most molecular bonds have their characteristics between these two extremes, between purely covalent and purely ionic bond
\end{itemize}
\end{minipage}

\xdef\loc{east}
\begin{tikzpicture}[remember picture,overlay] % anchor=south
    \node (A) [yshift=\figYshift,xshift=\figXshift, anchor=\loc] at (current page.\loc) {\includegraphics[page=1,width=4.8cm]{Figures_booklet/SOM_10_Bonding_figures/SOM_Bonding_Figure_11.png}}; %scale=\figscale. % 8cm max height
\end{tikzpicture}
\footnote{\HarrisCR[1-]}

\endofslide 
%\endofsection
\end{frame}
%%%%%%%%%%%%%%%

%%%%%%%%%%%%%%%
\begin{frame}
\frametitle{\texorpdfstring{\culine[\sectiontitleulinecolor]{\sectiontitle}}{\sectiontitle} }
\framesubtitle{Ioniset ja kovalenttiset sidokset}
\appear{}

\begin{itemize}[itemsep=6.5mm]
	\item Eli kahdenlaisia sidoksia \appear{(tai sidostyyppejä)} \appear{voi muodostua:} \appear{\CDAlert[blue]{ionisidoksia}}  \appear{ja  \CDAlert[red]{(polaarisia) kovalenttisia sidoksia}}\\[6.5mm]
\end{itemize}

% 6.5 cm = midway, 10 cm = 3/4 
\begin{minipage}{8.5cm}
\begin{itemize}[itemsep=6.5mm]

	\item Käytännössä valtaosassa molekyyleistä on sidoksia jotka ovat ominaisuuksiltaan jotain \Alert[fuchsia]{näiden ääriarvojen väliltä}\appear{, eivät ole puhtaasti ioni-/kovalenttisia sidoksia}
	
	\item Toisin sanoen \CDAlert[fuchsia]{binäärinen} jako kahden sidostyypin välillä \CDAlert[fuchsia]{on yksinkertaistus}\appear{, joka vain likimääräisesti kuvaa molekyylien välisiä sidoksia}
\end{itemize}
\end{minipage}

\xdef\loc{east}
\begin{tikzpicture}[remember picture,overlay] % anchor=south
    \node (A) [yshift=\figYshift,xshift=\figXshift, anchor=\loc] at (current page.\loc) {\includegraphics[page=1,width=4.8cm]{Figures_booklet/SOM_10_Bonding_figures/SOM_Bonding_Figure_11.png}}; %scale=\figscale. % 8cm max height
\end{tikzpicture}
\footnote{\HarrisCR[1-]}

\endofslide 
%\endofsection
\end{frame}
%%%%%%%%%%%%%%%

%%%%%%%%%%%%%%%
\begin{frame}
\frametitle{\texorpdfstring{\culine[\sectiontitleulinecolor]{\sectiontitle}}{\sectiontitle} }
\framesubtitle{Hiili, timantti ja sp$^3$-hybridisaatio}
\appear{}

% 6.5 cm = midway, 10 cm = 3/4 
\begin{minipage}{9.5cm}
	\begin{itemize}[itemsep=1.7mm]
	\item \CDAlert[fuchsia]{Hiilellä} on neljä $\mathrm{n}=2$ elektronia\appear{, kun hiiltä ympäröi korkean valenssin atomit}\appear{, hiili jakaa elektroninsa} 
	
	\item Spatiaalisesti tämä johtaa \appear{, neljään suuntaan osoittaviin }\appear{\Alert[fuchsia]{tetraedrimäisiin orbitaalikonfiguraatioihin}} %
	\appear{\xdef\loc{north east}%
\begin{tikzpicture}[remember picture,overlay]% anchor=south
    \node (A) [yshift=0.5*\figYshift,xshift=\figXshift, anchor=\loc] at (current page.\loc) {\includegraphics[page=1,width=4.1cm]{Figures_booklet/SOM_10_Bonding_figures/SOM_Bonding_Figure_12.png}};%scale=\figscale. % 8cm max height
\end{tikzpicture}\footnote{\HarrisCRshort}}

\item Nämä tilat tunnetaan \CDAlert[fuchsia]{$s p^{3}$-hybridiorbitaaleina} 

\item \CDAlert[fuchsia]{Timantti} on hyvä esimerkki%
\appear{\xdef\loc{south east}%
\begin{tikzpicture}[remember picture,overlay]% anchor=south
    \node (A) [yshift=-0.35*\figYshift,xshift=\figXshift, anchor=\loc] at (current page.\loc) {\includegraphics[page=1,width=4.1cm]{Figures_booklet/SOM_10_Bonding_figures/SOM_Bonding_Figure_13.png}};%scale=\figscale. % 8cm max height
\end{tikzpicture}%
}\appear{, jonka \Alert[fuchsia]{kide\-ra\-kenne on erittäin vahva}}\\[-1mm]
\[
{\small%
\begin{aligned}
\psi_{I} & \equal \appear{\psi_{2 s} } \minus  \appear{\psi_{2 p_{z}}} \\
\appear{\psi_{II}} &  \equal \appear{\psi_{2 s}}  \plus \frac{1}{3} \appear{\psi_{2 p_{z}}}  \minus \frac{8^{1/2}}{3} \appear{\psi_{2 p_{x}}} \\
\appear{\psi_{\mathrm{III}} } & \equal \appear{\psi_{2 s}+\Frac{1}{3} \psi_{2 p_{z}}} \appear{+\Frac{2^{1/2}}{3} \psi_{2 p_{x}}} \appear{-\Frac{6^{1/2}}{3} \psi_{2 p_{y}}} \\
\appear{\psi_{\mathrm{IV}}} & \equal  \appear{\psi_{2 s}}  \plus \frac{1}{3} \appear{\psi_{2 p_{z}}} \plus \frac{2^{1/2}}{3} \appear{\psi_{2 p_{x}}} \plus \frac{6^{1/2}}{3} \appear{\psi_{2 p_{y}}}
\end{aligned}
}
\]
\end{itemize}
\end{minipage}
%\vspace{3mm}




\endofslide 
%\endofsection
\end{frame}
%%%%%%%%%%%%%%%

%%%%%%%%%%%%%%%
\begin{frame}
\frametitle{\texorpdfstring{\culine[\sectiontitleulinecolor]{\sectiontitle}}{\sectiontitle} }
\framesubtitle{sp$^3$-hybridisaatio ja sähködipolimomentti}

\begin{itemize}[itemsep=3.5mm]
	\item \CDAlert[red]{Alkuainetaulukossa hiilen lähellä olevat alkuaineet muodostavat}  \CDAlert[red]{usein samanlaisia rakenteita} \appear{(esim. pii (Si)}\appear{, germanium (Ge)}\appear{, Si$_{1-x}$Ge$_{x}$ seokset}\appear{,...)}

\item Esimerkiksi typellä on viisi $\mathrm{n}=2$ electronia. \appear{Ammoniakkimolekyylin $ \textrm{((}\mathrm{NH}_{3} \textrm{))}$}%
\appear{\xdef\loc{south east}%
\begin{tikzpicture}[remember picture,overlay]% anchor=south
    \node (A) [yshift=-0.25*\figYshift,xshift=\figXshift, anchor=\loc] at (current page.\loc) {\includegraphics[page=1,width=7cm]{Figures_booklet/SOM_10_Bonding_figures/SOM_Bonding_Figure_15.png}};%scale=\figscale. % 8cm max height
\end{tikzpicture}\footnote{\HarrisCR}}%
\appear{ yksi viidestä elektronista ottaa 'kadoksissa' olevan vetyelektronin paikan.} \appear{Protonin ollessa myös kateissa orbitaalin tämän alueen sähkövaraus on negatiivinen} 
\implies{ Selittää ammoniakin \CDAlert[fuchsia]{sähködipolimomentin} \appear{mikä selittää esim.~sen \\ ominaisuudet \CDAlert[fuchsia]{liuottimena} }}
\end{itemize}

\vspace{2.5mm}
% 6.5 cm = midway, 10 cm = 3/4 
\begin{minipage}{6.7cm}
\begin{itemize}[itemsep=3.5mm]
	\item \CDAlert[blue]{Vedessä}\appear{, kaksi hapen ylimääräistä $\mathrm{n}=2$ elektronia korvaa vetyjen elektronit}
	\vspace{-0.5mm}
	
	\implies{ Kaksi orbitaalialuetta muodostuu \\ synnyttäen 
			myös \CDAlert[blue]{sähködipolimomentin} }
\end{itemize}
\end{minipage}



%\endofslide 
\endofsection
\end{frame}
%%%%%%%%%%%%%%%

